% ---
% id: meat-core/rem:repeated-executions
% kind: rem
% title: "One execution's security parameter and a reused key"
% summary: >-
%   One execution's parameter covers one run, and separate executions that
%   share no key add their parameters ([[cor:separate-executions]]). A reused
%   authentication key enters through ε_auth, charged twice, and the earlier
%   run's parameter does not bound its distance in general.
% user_statement: >-
%   The security parameter ε covers one execution: one run of n rounds and
%   its post-processing. It is not a per-round or per-bit figure, and it does
%   not shrink when you repeat the protocol. Several executions are together
%   secure with the sum of their ε ([[cor:separate-executions]]) under
%   [[A-SEPARATE]]. The sum does not cover parameters chosen from an earlier
%   run's data, such as the error-correction budget or the acceptance test.
%   The conditions of [[A-SEPARATE]] include: (i) their number and parameters
%   are all fixed before the first starts; (ii) they share no key material
%   and no device state; (iii) their keys stay unused until the last has
%   ended. Reusing key from an earlier run to authenticate the next is not
%   covered either. Such a key enters only through the next run's declared
%   authentication error ε_auth, counted twice, and the earlier run's ε does
%   not by itself bound how far that key is from an ideal one.
% prominence: standard
% status: accepted
% uses: [cor:composed-security, cor:separate-executions]
% mentions: [ass:auth, ass:separate-executions]
% assumes: []
% refs:
%   - "portmann2022security, Sec. II.E"
%   - "portmann2022security, Sec. II.F"
%   - "tupkary2025survey, Sec. 2.3, footnote 2"
%   - "tupkary2025survey, Sec. 2.3, Remark 2.1"
%   - "tupkary2026decoy, Sec. 3.3, Def. 3.8, Eqs. (11), (17)"
% provenance: newly-derived
% lean:
% ---

\begin{remark}\label{rem:repeated-executions}
The parameter $\epa+\eev+2\varepsilon_{\mathrm{auth}}$ of
Corollary~\ref{cor:composed-security}\,(iii) bounds one execution: one run
of $n$ rounds and its post-processing. It is not a rate per round or per
bit, and it does not shrink when executions are repeated.

A run whose pre-shared authentication key is not fresh, for instance part of
the key of an earlier execution, enters through
Assumption~\ref{ass:auth}\,(iii)--(v). If, before the run starts, the key
jointly with every other register and the adversary's information is within
trace distance $\delta$ of the state in which it is replaced by a uniform and
independent key, then running the scheme, the protocol and any
distinguisher on it is a channel applied to that joint state, together with
fresh systems independent of it; a channel does not increase the trace
distance, so the outputs differ by at most $\delta$ from those on the ideal
key, and by the triangle inequality the scheme realises the authenticated
classical channel within $\varepsilon_0+\delta$, where $\varepsilon_0$ is its
error on an ideal key. The declared $\varepsilon_{\mathrm{auth}}$ includes
$\delta$, and Corollary~\ref{cor:composed-security}\,(iii) charges it twice,
$2\delta$ in all. For a key taken from an earlier execution's output, that
execution's parameter does not bound $\delta$ in general. That parameter is
a sum over the pairs of key lengths
\cite[Sec.~3.3, Def.~3.8, Eqs.~(11), (17)]{tupkary2026decoy}: it bounds the
key's distance from an ideal one where both parties hold a key only
weighted by the probability of that event, which the adversary controls
\cite[Sec.~2.3, Remark~2.1]{tupkary2025survey}. Where one party holds a key
and the other none, as after a blocked final $\mathrm{accept}$ of the
\tref{authentication post-processing}, the ideal state of that definition
also leaves one party without a key, so the parameter does not bound the
probability of that event; there no shared key exists to reuse, and
$\delta$, the distance of the whole state on which the next run starts,
conditioned on its starting, counts that event where the next run starts
on it, and otherwise enters $\delta$ through the conditioning on its
start, which the adversary influences.

For executions that share no key material and no device state, and whose
keys stay unused until the last of them has reached its
\tref{end of an execution} (Assumption~\ref{ass:separate-executions}), the
parameters add:
Corollary~\ref{cor:separate-executions} bounds $N$ such executions, run
against one adversary in any order, by
$\varepsilon_1+\dots+\varepsilon_N$ for the joint ideal map, which is what
composition gives for constructive statements
\cite[Sec.~II.E]{portmann2022security}. Not covered are a key used before
the last execution has ended, since the ideal map applied at the end would
then replace a key whose use the adversary has already seen; parameters
chosen from an earlier execution's data, which would make each
$\varepsilon_i$ depend on the attack; and a chain in which a key of one
execution authenticates the next. For such a chain the hybrid would have
to replace a key by an ideal one at a moment both parties have ended the
earlier execution, which the authentication post-processing does not
provide, and the earlier parameter does not bound $\delta$ (the paragraph
above). That the literature expects such chains to add
\cite[Sec.~2.3, footnote~2]{tupkary2025survey} is cited as motivation
only, and a bound on the failure probability accumulated over every use
\cite[Sec.~II.F]{portmann2022security} lies outside what this contract
proves.
\end{remark}
